Die Wasserschicht hat das Volumen $V = A\,d$ und die Masse $m = \rho\,A\,d$ mit $\rho = \rhWO$ (FoTaBe S.~187). Mit $c = \cWO$ (FoTaBe S.~191) ist die nötige Wärme
\begin{align}
 Q &= \QZF \\
   &= \rhW\cdot\AZ\cdot\dS\cdot\cW\cdot\dT \\
   &= \QZ = \QZk \approx \result{\QZGP{0}{3}}
\end{align}
Beim Preis $p$ von \prkO\ Rappen pro Kilowattstunde kostet das
\begin{align}
 K &= \KosF = \QZk\cdot\prkO\,\frac{\text{Rp.}}{\si{\kilo\watt\hour}} \approx \result{\KosP{0}{2}\ \text{Mio.\ Franken}}
\end{align}
